\section{Whole rows, minimum dimension, and affine normalization}
\label{sec:whole-rows}
The differential bounds above allow one slack row to be split across
many factors. We first study the more restrictive \emph{whole-row}
model, in which the source bodies are divided into groups and each group
is assigned to one cone that factors all support rows of its bodies.
There the linear rank of the slack functions bounds the cone dimension,
the dimension of a contact face bounds the largest exposed-face
dimension, and one cone attains both bounds, without any regularity
assumption. We then prove two results that remain valid when rows split:
rigidity at minimum total dimension, and a dimension bound for
projections of a single normalized base of a product cone. Finally, we
characterize when one level set of a linear functional normalizes a lift
without changing its projection, so that every barrier on the lift
induces a barrier on the body with no larger parameter.

\subsection{Contact faces and exact grouping}
Let $C_a\subset\R^{n_a}$ be compact full-dimensional convex bodies with
$0\in\intr C_a$, for $a\in[k]$. Write
\[
 E_a=\operatorname{ext}C_a,\qquad P_a=\operatorname{ext}C_a^\circ,
 \qquad F_a(z)=\{x\in C_a:\ip{x}{z}=1\}.
\]
Define the maximum contact-face dimension and its codimension by
\[
 e_a=\max_{z\in P_a}\dim F_a(z),\qquad \lambda_a=n_a-e_a.
\]
Here face dimension is affine dimension, so $1\leq\lambda_a\leq n_a$.
The maximum exists because the nonempty set of possible dimensions is
a finite set of integers. For a proper cone $K$ define
\[
 f(K)=\max_{0\ne y\in K^*}\dim\spanop(K\cap y^\perp).
\]
Thus $f(K)$ is the largest dimension of the linear span of a proper
exposed face of $K$; the zero face is allowed.

\begin{theorem}[Whole-row dimension and face bounds]
\label{thm:whole-row-caps}
Fix a nonempty group $G\subseteq[k]$. Suppose arbitrary maps
\[
 A:\prod_{a\in G}E_a\longrightarrow K,\qquad
 B^a:P_a\longrightarrow K^*
\]
into a proper $m$-dimensional cone satisfy
\begin{equation}\label{eq:whole-row-slack}
 \ip{A(x)}{B^a(z)}=1-\ip{x_a}{z}
 \quad(a\in G,\ x\in\prod_{b\in G}E_b,\ z\in P_a).
\end{equation}
Put $N_G=\sum_{a\in G}n_a$. Then
\begin{equation}\label{eq:whole-row-bounds}
 m\geq N_G+1,\qquad f(K)\geq N_G+1-\min_{a\in G}\lambda_a.
\end{equation}
Both bounds are attained by the shared homogenization cone
\begin{equation}\label{eq:shared-homogenization}
 H_G=\{(t,(y_a)_{a\in G}):t\geq0,\ y_a\in tC_a\text{ for all }a\in G\}.
\end{equation}
No continuity or definability of the factor maps is required.
\end{theorem}
\begin{proof}
The sets $E_a$ and $P_a$ affinely span their respective spaces: their
convex hulls are $C_a$ and $C_a^\circ$. Consequently the functions
$1-\ip{x_a}{z}$, indexed by all $a,z$, span the constant and every
coordinate of $x$. These $N_G+1$ functions are linearly independent on
$\prod_{a\in G}E_a$: fix all coordinates but one and use its affine span,
successively for each block. By \eqref{eq:whole-row-slack}, their span
has dimension at most $m$. This proves the first bound.

Fix $a,z$ and restrict $x$ to the contact cylinder
\[
 (F_a(z)\cap E_a)\times\prod_{b\in G\setminus\{a\}}E_b.
\]
Its image under $A$ lies in $K\cap B^a(z)^\perp$. The vector $B^a(z)$
is nonzero, since its slack is not identically zero. The extreme points
of the exposed face $F_a(z)$ are precisely $F_a(z)\cap E_a$, and they
affinely span that face. Thus the restricted affine functions have rank
\[
 1+\dim F_a(z)+\sum_{b\in G\setminus\{a\}}n_b.
\]
They are restrictions of linear functionals of $A(x)$, so this quantity
cannot exceed $\dim\spanop(K\cap B^a(z)^\perp)\leq f(K)$.
Maximizing over $a,z$ proves the second bound.

The cone $H_G$ is proper and has dimension $N_G+1$. The maps
$A(x)=(1,x)$ and $B^a(z)=(1,0,\ldots,-z,\ldots,0)$ give the required
factorization. To compute all its exposed faces, write
$h_a(w)=\max_{x\in C_a}\ip{x}{w}$. Its dual is
\[
 H_G^*=\{(\alpha,(w_a)_a):\alpha\geq\sum_a h_a(-w_a)\}.
\]
An interior dual point exposes only zero. At a nonzero dual boundary
point let $J=\{a:w_a\ne0\}$, which is nonempty. Nonnegativity of each
support slack shows that the complementary face has dimension
\begin{equation}\label{eq:shared-face-dimension}
 1+\sum_{a\notin J}n_a+
 \sum_{a\in J}\dim\arg\min_{x\in C_a}\ip{x}{w_a}.
\end{equation}
Every face $F_a(z)$ with $z$ on the boundary of $C_a^\circ$ is an
intersection of faces $F_a(z')$ with $z'\in P_a$. Indeed, express $z$
as a finite convex combination of extreme points of the polar;
equality in its supporting inequality forces equality in each
inequality having positive weight. Each block $a\in J$ in
\eqref{eq:shared-face-dimension} therefore contributes at most
$n_a-\lambda_a$. A single block in $J$ with smallest $\lambda_a$, using
a normal attaining $e_a$, attains the claimed maximum.
\end{proof}

By a \emph{whole-row partition} we mean a partition $\mathcal G$ of
$[k]$ with one cone assigned to each group, factoring all of that group's
rows as in \eqref{eq:whole-row-slack}. In particular, no row is split
between different cones. Allowing a group's primal map to
depend on other source coordinates does not weaken the lower bound:
fix those coordinates and apply the theorem.

\begin{corollary}[Exact grouping under two caps]
\label{cor:whole-row-grouping}
Let $D\geq1$ and $F\geq0$ be integers, and allow all proper cones with
dimension at most $D$ and $f(K)\leq F$.
A group $G$ is feasible in the whole-row model if and only if
\begin{equation}\label{eq:whole-row-grouping}
 \sum_{a\in G}n_a\leq D-1,\qquad
 \sum_{a\in G}n_a-\min_{a\in G}\lambda_a\leq F-1.
\end{equation}
The minimum factor count is the minimum number of parts satisfying these
two inequalities. For $k$ identical bodies with dimension $n$ and contact
codimension $\lambda$, put
\[
 Q=\min\left\{\left\lfloor\frac{D-1}{n}\right\rfloor,
                 \left\lfloor\frac{F+\lambda-1}{n}\right\rfloor\right\}.
\]
If $Q\geq1$, the exact factor count is $\lceil k/Q\rceil$; if $Q=0$,
no whole-row partition is feasible.
\end{corollary}
\begin{proof}
Necessity is \eqref{eq:whole-row-bounds}, and $H_G$ proves sufficiency
for every feasible group. The identical-body inequalities reduce to
$|G|\leq Q$.
\end{proof}

For strictly convex bodies, $\lambda_a=n_a$. For polytopes,
$\lambda_a=1$, since some extreme polar normal exposes a facet. For
Euclidean balls of arbitrary dimensions, the theorem gives exactly
\[
 \dim H_G=1+\sum_{a\in G}n_a,\qquad
 f(H_G)=1+\sum_{a\in G}n_a-\min_{a\in G}n_a.
\]
Thus the face bound counts the full coordinate spaces of the other
bodies in the group, not only tangent directions detected by curvature.

\begin{proposition}[Complexity of heterogeneous grouping]
\label{prop:whole-row-grouping-hard}
For Euclidean balls with dimensions given as positive integers, deciding
whether the whole-row grouping problem has a partition with at most a
specified number of groups is strongly NP-complete. This holds already
when $F=D-1$.
\end{proposition}
\begin{proof}
Under $F=D-1$, the second inequality in
\eqref{eq:whole-row-grouping} follows from the first because
$\min_{a\in G}n_a\geq1$. Consider a 3-PARTITION instance with integers
$b_1,\ldots,b_{3q}$, $B/4<b_a<B/2$, and $\sum_a b_a=qB$
\cite[Problem SP15]{GareyJohnson1979}. Set $n_a=b_a$, $D=B+1$,
$F=B$, and ask for at most $q$ groups. Their total dimension is $qB$,
so every feasible group must have dimension exactly $B$ and there must
be exactly $q$ groups. The size inequalities force exactly three items
in each group. This is precisely the original 3-PARTITION question.
The reduction preserves polynomial bounds on all integers, so on the
polynomially bounded instances used for strong NP-completeness even
explicit ball descriptions have polynomial size. A partition is a
polynomially checkable certificate for the two integer inequalities,
which also proves membership in NP. The hardness comes from the
classical packing problem, not from testing a fixed group.
\end{proof}

\subsection{Rigidity at the absolute dimension minimum}
The next result applies to arbitrary compact bodies and allows every
row to split. A lift of minimum total dimension is, up to a linear
isomorphism, the homogenization cone of the body, so it cannot improve
any linear-isomorphism invariant.

\begin{theorem}[Minimum-dimensional rigidity]
\label{thm:minimal-dimension-rigidity}
Let $C\subset\R^N$, $N\geq1$, be a compact full-dimensional convex body.
Every exact lift over a proper cone of total dimension $M$, allowing
affine slices, affine projections, and free variables, has $M\geq N+1$.
At equality, after the reduction of Lemma~\ref{lem:reduction}, an
invertible linear map carries the cone to
\[
 \operatorname{hom}(C)=\{(t,tx):t\geq0,\ x\in C\}
\]
and its feasible slice to $\{1\}\times C$. The reduced projection is
an affine isomorphism on that slice.

The same cone-isomorphism conclusion holds for any factorization of
the full slack operator through a proper cone of dimension $N+1$,
without regularity assumptions on its factor maps. In particular, it
holds when that cone is a product and rows split across its factors.
\end{theorem}
\begin{proof}
First consider an exact lift. Eliminate free variables and reduce to the
minimal face, of dimension $M'\leq M$, with a Slater affine slice $A$.
Its affine output has full rank on the direction space of $A$, so
$\dim A\geq N$. If $M'\leq N$, these facts force $A$ to be the entire
cone space. Its image is an affine translate of a cone, which cannot
be a full-dimensional bounded body. Thus $M'\geq N+1$.

When $M=N+1$, no face reduction is possible, and $A$ must be an
$N$-dimensional affine hyperplane: the full-space case gives the same
contradiction. The output map on $A$ is an affine isomorphism, so
$S=K\cap A$ is compact. The hyperplane does not contain zero, again
because otherwise its feasible set and output would be translated
cones. Write $A=\{z:\ell(z)=1\}$. Fix $w\in A\cap\intr K$.
If nonzero $v\in K$ satisfies $\ell(v)=0$, then $w+tv\in S$ for all
$t\geq0$, contradicting compactness. If $\ell(v)<0$, the nonzero
vector $v-\ell(v)w\in K\cap\ker\ell$ gives the same contradiction;
pointedness ensures this vector is nonzero. Therefore $\ell$ is
strictly positive on $K\setminus\{0\}$, and $S$ is a compact base.
The affine output on $A$ extends to a linear map $\widetilde\pi$ on
the cone space because $A$ is a hyperplane not containing zero. The map
$Tz=(\ell(z),\widetilde\pi(z))$ is invertible: its kernel lies in the
direction space of $A$, where the output is injective. Writing each
nonzero $z\in K$ uniquely as $z=\ell(z)s$, $s\in S$, proves both claims.

For the factorization assertion, translate $C$ so $0\in\intr C$.
The canonical vectors $U(x)=(1,x)$ for $x\in\operatorname{ext}C$ and
$V(z)=(1,-z)$ for $z\in\operatorname{ext}C^\circ$ each span
$\R^{N+1}$. The slack has rank $N+1$, by the affine-function argument
in Theorem~\ref{thm:whole-row-caps}. If $A(x),B(z)$ give another
factorization in that dimension, their images span as well. A linear
relation among the $U(x)$ is equivalent, by pairing with the spanning
dual family, to the same relation among the $A(x)$. Hence an invertible
$T$ satisfies $A=TU$ and $B=T^{-\mathsf T}V$. Taking conic hulls gives
\[
 T\operatorname{hom}(C)\subseteq K,\qquad
 T^{-\mathsf T}\operatorname{hom}(C)^*\subseteq K^*.
\]
The extreme points of the body and its polar generate their respective
compact bases, so these are the full homogenization cones. Dualizing
the second inclusion reverses the first, and proves equality.
\end{proof}

The rows in Theorem~\ref{thm:whole-row-caps} are exactly the slack rows
of $\prod_{a\in G}C_a$ indexed by the extreme points of its polar.
Indeed, its polar is
$\{(z_a):\sum_a h_a(z_a)\leq1\}$. A nonzero extreme point must have
one supported coordinate block, since two positive support-function
contributions give a nontrivial convex decomposition into single-block
points. Within that block it must be an extreme point of $C_a^\circ$.
Conversely each such single-block point is extreme. Thus equality
$m=N_G+1$ in that theorem forces $K$ to be linearly isomorphic to $H_G$.
For a fixed partition into $g$ groups, every whole-row formulation of
total dimension $\sum_a n_a+g$ therefore uses precisely those shared
homogenization cones up to linear changes of coordinates. Barrier
parameters are preserved by such changes. In particular, the exact
shared spectral-cone values of
Theorem~\ref{thm:shared-spectral-parameter}, including its Euclidean
vector-ball case, apply to every such minimum-dimensional formulation.

\begin{corollary}[Connected extreme sets under a strict factor cap]
\label{cor:connected-extreme-gap}
Suppose $\operatorname{ext}C$ is connected. A full slack factorization
through $K=\prod_{i=1}^L K_i$, with every $K_i$ nonzero and proper,
can have total dimension $N+1$ only when $L=1$. Consequently, if
$\dim K_i\leq d<N+1$, then
\begin{equation}\label{eq:connected-extreme-gap}
 M\geq N+2,\qquad
 L\geq\max\left\{2,\left\lceil\frac{N+2}{d}\right\rceil\right\}.
\end{equation}
These conclusions also hold for exact lifts after the reduction of
Lemma~\ref{lem:reduction}, and allow arbitrary row splitting.
\end{corollary}
\begin{proof}
Normalize extreme rays by any strictly positive linear functional.
For a product of $L$ nonzero proper cones, the extreme rays are supported
in exactly one factor; the $L$ factor-supported pieces are nonempty
and both open and closed in the normalized extreme-ray space. For
$\operatorname{hom}(C)$ that space is homeomorphic to
$\operatorname{ext}C$. Theorem~\ref{thm:minimal-dimension-rigidity}
therefore excludes $L>1$ at equality. A single factor under the strict
cap cannot supply dimension $N+1$. Integrality gives $M\geq N+2$,
and $M\leq Ld$ gives the count bound. For an unreduced exact lift
with $M=N+1$, the theorem already ensures that no reduction occurs.
\end{proof}

The corollary excludes only $M=N+1$; it does not force one extra
coordinate per factor. The balance slices in
Theorem~\ref{thm:balance-slices} have connected extreme sets and $M=N+2$
for arbitrarily many essential factors, so $M\geq N+L$ fails in general.

\subsection{What one normalization equation does imply}
\begin{proposition}[An additive bound for one normalized base]
\label{prop:normalized-base-tax}
Let $K=\prod_{i=1}^L K_i$, where $K_i$ is nonzero and proper of
dimension $m_i$, and choose $\ell_i\in\intr K_i^*$. Let
\[
 B=\{z\in K:\sum_i\ell_i(z_i)=1\},\qquad C=\pi(B),
\]
where $\pi$ is affine and $\dim\operatorname{aff}C=N$.
If $\operatorname{ext}C$ is connected, then
$\sum_i m_i\geq N+L$. No irredundancy assumption is needed.
\end{proposition}
\begin{proof}
Let $B_i$ be the compact normalized base of factor $i$, embedded in
the product with all other coordinates zero, and put $C_i=\pi(B_i)$.
Then $C=\operatorname{conv}\bigcup_i C_i$, and
$\dim\operatorname{aff}C_i\leq m_i-1$. Every extreme point of $C$
belongs to some $C_i$, because any finite convex representation of an
extreme point uses only that point. The nonempty sets
$E_i=C_i\cap\operatorname{ext}C$ form a finite relatively closed cover
of the connected extreme set. Their intersection graph is connected:
a graph separation would give a separation into two disjoint nonempty
relatively closed finite unions. Along a spanning tree, the affine
hull of each newly added $E_i$ intersects the previous affine hull.
The dimension of their combined affine hull is therefore at most
the sum of their dimensions. With $I=\{i:E_i\ne\varnothing\}$ this gives
\[
 N=\dim\operatorname{aff}(\operatorname{ext}C)
 \leq\sum_{i\in I}\dim\operatorname{aff}E_i
 \leq\sum_{i\in I}(m_i-1)
 \leq\sum_{i=1}^L(m_i-1).
\]
The final inequality uses $m_i\geq1$. Rearrangement proves the result.
\end{proof}

A second affine equation can create extreme points with nonzero
coordinates in several factors. The sets $C_i$ in the proof then need not
cover the extreme points of the smaller slice, so the proposition does
not contradict the balance-slice examples.

\subsection{Recession normalization and induced barriers}
The projection lemma in Section~\ref{sec:concrete-barriers} requires
bounded fibers. The next criterion states exactly when one level set of
a linear functional provides them without changing the projection.
All interiors in the statement are
relative to affine hulls.

\begin{proposition}[A criterion for barrier-preserving normalization]
\label{prop:recession-normalization}
Let $D$ be a nonempty closed convex set containing no affine line,
and let an affine map $\pi$ satisfy $\pi(D)=C$, where $C$ is a compact
convex body. Write $R=\operatorname{rec}D$ for its recession cone.
If $R\ne\{0\}$, choose a linear functional $a$ strictly positive on
$R\setminus\{0\}$, and define
\[
 m_a(x)=\min\{a(z):z\in D,\ \pi(z)=x\}.
\]
Every minimum is finite and attained. For any real $M$,
\begin{equation}\label{eq:recession-normalization}
 \pi(D\cap\{a=M\})=C
 \quad\Longleftrightarrow\quad M\geq\sup_{x\in C}m_a(x).
\end{equation}
Whenever this holds, the level section is compact. If the supremum is
finite, every self-concordant barrier on $\ri D$ induces a barrier on
$\intr C$ with no larger parameter. If $R=\{0\}$, the same parameter
conclusion holds without taking a level section.
\end{proposition}
\begin{proof}
Compactness of $C$ implies $R\subseteq\ker\pi_0$, where $\pi_0$ is
the linear part of $\pi$. The cone $R$ is pointed because $D$ contains
no line, so a strictly positive $a$ exists. Every nonempty sublevel
$D\cap\{a\leq b\}$ is bounded. Otherwise normalize differences of an
unbounded sequence from a fixed point of $D$ and pass to a unit limit
$r$. Closedness and convexity imply $r\in R$, while boundedness above
of $a$ on the sequence gives $a(r)\leq0$, a contradiction. Each fiber
is closed; minimizing $a$ on a compact nonempty sublevel proves
attainment and finiteness of $m_a(x)$.

Fix $0\ne r_0\in R$. Adding a nonnegative multiple of $r_0$ to a
minimizer in a fiber preserves that fiber and reaches every larger
$a$-value. This proves \eqref{eq:recession-normalization}. The level
section is closed and bounded by the same sublevel argument.
If the supremum is finite, choose $M$ greater than it and greater than
$a(z^0)$ for some $z^0\in\ri D$. Adding a suitable multiple of $r_0$
to $z^0$ gives a point of the level in $\ri D$. Indeed, a relative
ball around $z^0$ remains in $D$ after translating by a recession
direction. Thus the affine restriction of the original barrier is a
barrier on the relative interior of the compact level section, with
the same parameter bound. Lemma~\ref{lem:bounded-fiber-barrier} then
gives a barrier on $C$. If $R=\{0\}$, the closed convex set $D$ is
bounded by the same normalized-sequence argument, and that lemma
applies directly.
\end{proof}

\begin{corollary}[Arbitrary line-free lifts of a simple-vertex polytope]
\label{cor:polytope-lift-barrier}
Let $C$ be an $N$-dimensional compact polytope with a simple vertex.
Every self-concordant barrier on the relative interior of every closed
line-free affine lift of $C$ has parameter at least $N$, even if its
fibers are unbounded. The minimum over all such lifts and barriers is $N$.
\end{corollary}
\begin{proof}
Choose a lift $z_v\in D$ of each vertex $v$ of $C$. For a convex
combination $x=\sum_v\theta_vv$, its lifted combination belongs to $D$
and therefore
\[
 m_a(x)\leq\sum_v\theta_va(z_v)\leq\max_v a(z_v).
\]
By Proposition~\ref{prop:recession-normalization}, the barrier induces a
barrier on $C$ with no larger parameter. At a simple vertex, the
classical independent-facet lower bound gives $\nu\geq N$
\cite[Proposition~2.3.6]{NN1994}. The universal barrier on $C$ attains
the bound $N$ \cite[Theorem~2]{LeeYue2021}; it may also be used on
the canonical homogenization slice, which is affinely isomorphic to $C$.
\end{proof}

The simple-vertex hypothesis is essential to this argument. Choosing
$N$ independent normals at a nonsimple vertex does not reproduce its
tangent cone. At $e_1$ of the three-dimensional cross-polytope, the four
active normals are $(1,\pm1,\pm1)$. The three with sign pairs
$(+,+),(+,-),(-,+)$ are independent, but $(-1,-1,-1)$ satisfies their
three strict tangent inequalities and violates the fourth. The cited
simple-vertex lower bound cannot be applied after deleting that facet.

\begin{example}[Compact projection need not admit affine normalization]
\label{ex:no-affine-normalization}
Let $D$ be the positive-semidefinite lift of the disk in
Example~\ref{ex:escaping-disk}, with coordinates $(x_1,x_2,u,z)$.
Its recession cone is exactly the positive $z$-ray. Indeed, compactness
of the projected disk makes every recession direction have zero
$x$-coordinates. The second block then requires its $u$-coordinate
to be nonnegative, while the last homogeneous block
$\left(\begin{smallmatrix}0&\dot u\\\dot u&\dot z\end{smallmatrix}\right)$
forces $\dot u=0$ and $\dot z\geq0$. Conversely every such direction
is a recession direction. If an affine subspace removed this ray from
a nonempty section of $D$, that closed convex section would have zero
recession cone and hence be bounded. It could not project onto the
entire disk, since every completion of the boundary points approaching
$(1,0)$ in Example~\ref{ex:escaping-disk} has $z\to\infty$.
Thus no affine restriction that preserves the projection removes the
recession ray, even for this semialgebraic Slater lift over three
$2\times2$ positive-semidefinite cones. The example rules out only
normalization by an affine restriction; it does not show that projecting
a general lift with unbounded fibers can increase the barrier parameter.
\end{example}

The slack-factorization framework used in this section is classical
\cite{GPT2013}. The minimum-dimension rigidity theorem is an elementary
consequence of it, and the normalized-base dimension bound uses only a
finite connected cover. The barrier bounds in
Corollary~\ref{cor:polytope-lift-barrier} are also classical; the
corollary adds that the lower bound holds on any given closed line-free
lift, by Proposition~\ref{prop:recession-normalization} and
Lemma~\ref{lem:bounded-fiber-barrier}. The whole-row grouping formulas
apply only to the whole-row model; they do not claim optimality among
lifts in which rows may split.
